\chapter{Rediseno UX/UI con paleta UAM-C}
\label{ch:rediseno}

\section{Objetivo}

Mejorar el acabado visual de la primera version del cliente:

\begin{itemize}[leftmargin=*]
  \item Definir una paleta inspirada en el naranja institucional de la UAM Cuajimalpa.
  \item Centralizar los estilos en un unico \texttt{ThemeData}.
  \item Eliminar el problema de iconos de Material que se renderizaban como glifos de fuentes CJK.
  \item Redisenar las pantallas de login y registro con el nuevo estilo.
\end{itemize}

\section{Contexto}

Al llegar a la Sesion~\ref{ch:sesion6b}, la app funcionaba pero se veia generica. Ademas, los iconos de Material Design aparecian como caracteres raros en la barra de navegacion inferior, en el campo de busqueda y en el boton de logout.

\section{Desarrollo}

\subsection{Diagnostico del problema de iconos}

Los glifos ``raros'' no eran caracteres chinos insertados en el codigo. Eran iconos de Material Design que no podian cargarse porque faltaba la directiva \texttt{uses-material-design: true} en el \texttt{pubspec.yaml}. Al no encontrarlos, Flutter caia en fuentes de respaldo del sistema (CJK), lo que producia los glifos incorrectos.

\subsection{Correccion del \texttt{pubspec.yaml}}

\begin{lstlisting}[language=bash]
name: notiuam_mobile
description: "Plataforma distribuida de comunicacion y alertas comunitarias."
publish_to: 'none'
version: 1.0.0+1

environment:
  sdk: ^3.13.4

dependencies:
  flutter:
    sdk: flutter
  cupertino_icons: ^1.0.8
  dio: ^5.7.0
  flutter_secure_storage: ^9.2.2
  provider: ^6.1.2
  intl: ^0.19.0

dev_dependencies:
  flutter_test:
    sdk: flutter
  flutter_lints: ^5.0.0

flutter:
  uses-material-design: true
\end{lstlisting}

\subsection{Paleta institucional}

Se definio una paleta que parte del naranja institucional:

\begin{table}[htbp]
\centering
\begin{tabular}{l l l}
\toprule
Rol & Hex & Uso \\
\midrule
\texttt{orange} & \#E85D04 & Botones, AppBar \\
\texttt{orangeDark} & \#B84700 & Estado pressed \\
\texttt{orangeLight} & \#F48C06 & Chips, badges \\
\texttt{charcoal} & \#2B2B2B & Texto principal \\
\texttt{cream} & \#FAFAF7 & Fondo \\
\texttt{success} & \#2A9D8F & Acciones positivas \\
\texttt{danger} & \#D62828 & Reportes, eliminar \\
\texttt{info} & \#457B9D & Informacion \\
\bottomrule
\end{tabular}
\caption{Paleta UAM-C utilizada en el cliente Flutter}
\end{table}

\subsection{Centralizacion del tema}

Se creo \texttt{lib/core/theme.dart} con dos clases: \texttt{NotiuamColors} (constantes) y \texttt{NotiuamTheme} (constructor del \texttt{ThemeData}).

\begin{lstlisting}[language=Java]
class NotiuamTheme {
  static ThemeData light() {
    final base = ThemeData(
      useMaterial3: true,
      colorScheme: ColorScheme.fromSeed(
        seedColor: NotiuamColors.orange,
        primary: NotiuamColors.orange,
        brightness: Brightness.light,
      ),
    );

    return base.copyWith(
      scaffoldBackgroundColor: NotiuamColors.cream,
      appBarTheme: const AppBarTheme(
        backgroundColor: NotiuamColors.orange,
        foregroundColor: Colors.white,
      ),
      filledButtonTheme: ...,
      outlinedButtonTheme: ...,
      inputDecorationTheme: ...,
      chipTheme: ...,
      navigationBarTheme: ...,
    );
  }
}
\end{lstlisting}

\subsection{Rediseno de pantallas}

\subsubsection{Login}

Se agrego un bloque de logo cuadrado con la letra N, subtitulo descriptivo, inputs con iconos, caja de error con fondo rosado y boton de altura fija con estado de carga.

\subsubsection{Registro}

Se anadio un encabezado con titulo y subtitulo, campos con iconos y limitacion de ancho a 420 px con \texttt{ConstrainedBox}.

\subsection{Cambios a archivos previos}

\begin{itemize}[leftmargin=*]
  \item \texttt{pubspec.yaml}: se reescribio agregando \texttt{uses-material-design: true}.
  \item \texttt{main.dart}: se cambio \texttt{theme: ThemeData(...)} por \texttt{theme: NotiuamTheme.light()}.
  \item \texttt{login\_screen.dart} y \texttt{register\_screen.dart}: rediseno completo.
  \item \texttt{explore\_screen.dart}: se ajustaron los colores del metodo \texttt{\_categoryIcon}.
\end{itemize}

\section{Verificacion}

\begin{itemize}[leftmargin=*]
  \item Los iconos de la barra inferior se ven como brujula, carpeta, marcador y persona.
  \item El AppBar y los botones principales usan el naranja institucional.
  \item Las tarjetas usan fondo blanco sobre crema.
  \item Los campos de texto tienen borde de foco naranja.
\end{itemize}

\section{Resultado}

\begin{exito}[Estado al cierre]
\begin{itemize}[leftmargin=*]
  \item Tema visual coherente en toda la app.
  \item Iconos de Material correctos.
  \item Base para las pantallas de administracion y moderacion.
\end{itemize}
\end{exito}